\section*{Capítulo \ref{ch:b2:wilson-cycle} --- Macizos antiguos y ciclo de Wilson}

\begin{solution}{exo:b2:wilson-cycle:1}
(a) $600 - 400 = \qty{200}{Myr}$, desde el rifting de los márgenes hasta la última colisión.
(b) $480 - 400 = \qty{80}{Myr}$. (c) La orogenia caledoniana es un conjunto de fases repartidas a lo largo de
unos 200 millones de años, y cada fase empezó en momentos distintos en lugares distintos: la palabra
sitúa un océano y su cordillera, no un momento. Una edad se da con la escala estratigráfica o en
millones de años.
\end{solution}

\begin{solution}{exo:b2:wilson-cycle:2}
(c) el rift de África oriental, rifting; (a) el mar Rojo, un océano joven; (d) el Atlántico, un océano
amplio con márgenes pasivos; (e) el margen pacífico de Sudamérica, subducción y
cierre de un océano; (b) el Himalaya, colisión; (f) Bretaña, una cordillera arrasada, tras lo cual
puede volver a empezar el rifting.
\end{solution}

\begin{solution}{exo:b2:wilson-cycle:3}
(a) Después del depósito de las capas del Silúrico inferior, es decir, después de \qty{443.1}{Ma}, y antes del
depósito de las capas del Devónico Superior, es decir, antes de \qty{358.9}{Ma}, el final del Devónico.
(b) Menos de $443.1 - 358.9 = \qty{84.2}{Myr}$; de hecho menos, ya que las capas inferiores son más jóvenes
que la base del Silúrico y las superiores más antiguas que su techo. (c) Depósito de las
turbiditas grises en el océano; plegamiento y esquistosidad en las colisiones caledonianas; levantamiento y
erosión hasta una superficie terrestre irregular; depósito de las brechas y conglomerados rojos; suave
basculamiento en la orogenia varisca; la erosión actual, que talla el acantilado.
\end{solution}

\begin{solution}{exo:b2:wilson-cycle:4}
(a) Ambas cosas: un bloque de corteza continental separado de los continentes por océanos mientras
atravesaba el Iapetus, y, una vez acoplado, un bloque limitado por fallas cuya historia difiere de la de sus
vecinos. (b) Ninguna: un continente, no un bloque añadido a uno. (c) Un \omterm{def:b2:wilson-cycle:terrane}{terrane}, no un
\omterm{def:b2:wilson-cycle:terrane}{microcontinente}: su corteza es la de un arco, no un fragmento de continente. (d) Ninguna: es
el límite entre \omterm{def:b2:wilson-cycle:terrane}{terranes}, la traza del océano cerrado.
\end{solution}

\begin{solution}{exo:b2:wilson-cycle:5}
(a) $3000 - 2600 = 400$, $2600 - 1500 = 1100$, $1500 - 875 = 625$, $875 - 600 = 275$, $600 -
262.5 = 337.5$; media $2737.5/5 = \qty{547.5}{Myr}$. (b) $1100/275 = 4.0$. (c) $(1500 -
262.5)/3 = \qty{412.5}{Myr}$. La media cambia en una cuarta parte al descartar los puntos dudosos
y las separaciones varían en un factor cuatro: una recurrencia aproximada, no un periodo.
\end{solution}

\begin{solution}{exo:b2:wilson-cycle:6}
(a) $3000/(70 - 25) = \qty{66.7}{Myr}$. (b) $3000/70 = \qty{42.9}{Myr}$. (c) Una vez subducida
la dorsal, el fondo oceánico que queda forma parte de la placa opuesta, que lleva el continente
opuesto: la velocidad de cierre es la velocidad de esa placa respecto de la fosa.
\end{solution}

\begin{solution}{exo:b2:wilson-cycle:7}
(a) $700/8 = \qty{87.5}{km}$; $700/30 = \qty{23.3}{km}$. (b) La cordillera varisca: con
\qty{30}{\celsius} por kilómetro funde la corteza media; con \qty{8}{\celsius} por kilómetro
la temperatura de fusión queda por debajo de la base incluso de una corteza engrosada. (c) El gradiente bajo
describe rocas arrastradas rápidamente hacia abajo en una zona de subducción, antes de que pudieran calentarse. Las rocas que
permanecen en profundidad en una corteza engrosada se calientan hacia una geoterma normal a lo largo de millones de años, y pueden
entonces alcanzar la fusión a menor profundidad.
\end{solution}

\begin{solution}{exo:b2:wilson-cycle:8}
(a) $\tau\ln(3/0.3) = 26.3\times2.303 = \qty{60.6}{Myr}$. (b) $290 - 60.6 = \qty{229}{Ma}$, en
el Triásico Superior. (c) $\tau = (290 - 251.9)/\ln10 = \qty{16.5}{Myr}$. (d) El colapso
extensional de la corteza engrosada, que la adelgaza y rebaja directamente su superficie; también una
tasa de erosión mayor, un clima más húmedo o rocas más blandas, o una altitud inicial menor.
\end{solution}

\begin{solution}{exo:b2:wilson-cycle:9}
(a) La sutura, traza del Iapetus, está en tierra firme en Gran Bretaña e Irlanda, al este del
Atlántico; corteza de Laurentia, del lado opuesto del Iapetus, se encuentra cerca de Rockall en el
lado europeo. El nuevo rift cortó al oeste de la antigua sutura: se abrió cerca de ella, no a lo largo de ella.
(b) Al norte de la sutura, el zócalo y la cobertera del margen laurentiano; al sur, los de
Avalonia, con historias distintas y, a comienzos del Paleozoico, faunas distintas. (c) El rifting
llegó al cabo de 80 millones de años, más que los 30 a 50 tras los cuales una zona enfriada puede haberse vuelto
resistente; y sin embargo el Atlántico Central sí siguió los Apalaches. El calor de una cordillera se desvanece; sus
fallas, fábricas y composiciones permanecen y aún pueden guiar un rift.
\end{solution}

\begin{solution}{exo:b2:wilson-cycle:10}
(a) $70 - 40/2 = \qty{50}{mm}$ al año. (b) $\qty{1500}{km}/\qty{50}{km.Myr^{-1}} =
\qty{30}{Myr}$. (c) El océano se estrecha a $70 - 40 = \qty{30}{mm}$ al año: $3000 - 30\times30 =
\qty{2100}{km}$. (d) Llega entonces a la fosa litosfera joven y caliente, y ya no hay ninguna placa vieja y fría que
se deshidrate bajo el arco: el magmatismo de arco cambia o cesa.
\end{solution}

\begin{solution}{exo:b2:wilson-cycle:11}
(a) $h = \Delta T(\rho_m - \rho_c)/\rho_m = 25\times600/3300 = \qty{4.5}{km}$. (b) Todos los
\qty{25}{km} adicionales, 5.5 veces la altura de la cordillera. (c) $2700\times9.8\times25\,000 =
\qty{6.6e8}{Pa}$, unos \qty{0.66}{GPa}. (d) A unos \qty{0.7}{GPa} y a las temperaturas de
la raíz de una corteza engrosada, varios cientos de grados, la facies de anfibolitas o la de
granulitas: los gneises y las \omterm{def:b2:crustal-melting:migmatite}{migmatitas} de los macizos antiguos.
\end{solution}

\begin{solution}{exo:b2:wilson-cycle:12}
(a) $S = 30\times10 = \qty{300}{km}$; $W_0 = 300 + 300 = \qty{600}{km}$; $T = 35\times600/300 =
\qty{70}{km}$. (b) $h = 35\times600/3300 = \qty{6.36}{km}$; $r = 4.5h = \qty{28.6}{km}$.
(c) $\Delta F = 2700\times9.8\times6364\times(35\,000 + 3182 + 14\,318) =
\qty{8.8e12}{N.m^{-1}}$, 1.4 veces el del Tíbet. (d) Una cordillera así pesa más que el Tíbet, que ya
se extiende: \omterm{def:b2:himalaya-andes:collapse}{colapso gravitacional}, por extensión de sus regiones altas y cabalgamiento en sus bordes.
\end{solution}

\begin{solution}{pb:b2:wilson-cycle:1}
\textbf{1.} $720/(540 - 500) = \qty{18}{km}$ por millón de años, \qty{18}{mm} al año.
\textbf{2.} \qty{36}{mm} al año: una dorsal lenta (de 20 a \qty{55}{mm} al año).
\textbf{3.} $36\times(580 - 540) = \qty{1440}{km}$.
\textbf{4.} $36\times(580 - 500) = \qty{2880}{km}$.
\textbf{5.} Expansión simétrica a velocidad constante desde anchura cero a \qty{580}{Ma}, y ninguna
subducción antes de \qty{500}{Ma}.
\textbf{6.} \qty{80}{Myr}, en los dos bordes del océano, junto a los márgenes continentales. Es la
litosfera oceánica más fría, más gruesa y más densa del océano, la más dispuesta a hundirse.
\textbf{7.} $64 - 36 = \qty{28}{mm}$ al año.
\textbf{8.} $2880/28 = \qty{102.9}{Myr}$ después de \qty{500}{Ma}: a \qty{397.1}{Ma}.
\textbf{9.} La dorsal, a \qty{1440}{km} de la fosa, se acerca a ella a $64 - 18 =
\qty{46}{mm}$ al año: al cabo de $1440/46 = \qty{31.3}{Myr}$, a \qty{468.7}{Ma}.
\textbf{10.} $2880 - 28\times31.3 = \qty{2004}{km}$.
\textbf{11.} $2004/40 = \qty{50.1}{Myr}$ después: a $468.7 - 50.1 = \qty{418.6}{Ma}$.
\textbf{12.} Dentro de los 425 a \qty{400}{Ma} de la colisión escandiana.
\textbf{13.} $40\times8 = \qty{320}{km}$.
\textbf{14.} $W_0 = 400 + 320 = \qty{720}{km}$; $T = 35\times720/400 = \qty{63.0}{km}$.
\textbf{15.} $h = 28\times600/3300 = \qty{5.09}{km}$; $r = 4.5\times5.09 = \qty{22.9}{km}$.
\textbf{16.} $\Delta F = 2700\times9.8\times5091\times(35\,000 + 2545 + 11\,455) =
\qty{6.60e12}{N.m^{-1}}$, cerca de un \qty{8}{\percent} más que el del Tíbet.
\textbf{17.} Dejar de elevarse y extenderse: extensión y fallas normales en su alto interior,
cabalgamientos en sus bordes, como hace hoy el Tíbet.
\textbf{18.} $f = 600/3300 = 0.182$; $k = 2.303/11 = 0.209$ por millón de años; $\tau = 1/(kf)
= \qty{26.3}{Myr}$.
\textbf{19.} $3.0\,e^{-30/26.3} = \qty{0.96}{km}$.
\textbf{20.} $26.3\ln(3.0/0.5) = \qty{47.1}{Myr}$ después de \qty{410}{Ma}: a \qty{362.9}{Ma}.
\textbf{21.} La corteza adicional es $h/f$: $3.0/0.182 = \qty{16.5}{km}$ baja a $0.5/0.182 =
\qty{2.75}{km}$, de modo que se retiran \qty{13.75}{km} de roca.
\textbf{22.} $2700\times9.8\times13\,750 = \qty{3.6e8}{Pa}$, unos \qty{0.36}{GPa}.
\textbf{23.} Sí: en el Devónico Superior la cordillera modelo ha bajado a cerca de \qty{1}{km} de altitud
media y sigue bajando, con relieve entre sus valles y crestas; las gravas rojas pueden
acumularse en sus zonas bajas mucho antes de \qty{358.9}{Ma}, como en Siccar Point, donde descansan sobre una superficie
con relieve.
\textbf{24.} \textbf{\qty[text-series-to-math=true]{2880}{km}}, a \qty{500}{Ma}, cuando
empezó la subducción.
\end{solution}
